\documentclass[a4paper,11pt]{article}
\usepackage[utf8]{inputenc}
\usepackage[T1]{fontenc}
\usepackage[english]{babel}
\usepackage[left=2.5cm,right=2.5cm,top=2cm,bottom=2cm]{geometry}
\usepackage{hyperref}
\usepackage{xcolor}
\usepackage{microtype}
\usepackage{lmodern}

\definecolor{primary}{RGB}{0, 85, 164}
\hypersetup{colorlinks=true, urlcolor=primary}

\begin{document}
\noindent \textbf{Lo\"ic Aron MBASSI EWOLO} \\
ENSEA -- Double Dipl\^ome Ing\'enieur \\
\href{mailto:loicaron.mbassiewolo@ensea.fr}{loicaron.mbassiewolo@ensea.fr} \quad (+33) 7 59 52 33 98

\vspace{1em}\noindent Cergy, le 21 mars 2026
\vspace{1em}\noindent \textbf{Objet : Candidature -- Alternance Data Scientist R\&D}
\vspace{1.5em}\noindent Madame, Monsieur,

\vspace{0.8em}\noindent Ing\'enieur ENSEA/Polytechnique Yaounde sp\'ecialis\'e en math\'ematiques appliqu\'ees et IA, je suis particuli\`erement attir\'e par la mission data scientist R\&D que propose RTE pour le syst\`eme Scoop (planification de tourn\'ees r\'eseau). J'ai d\'evelopp\'e au cours du projet Urban Waste Management (2024) un syst\`eme int\'egr\'e conjuguant mod\`eles ML de pr\'ediction (ARIMA, Random Forest) pour anticiper les besoins de collecte et algorithmes d'optimisation combinatoire (TSP, heuristiques) pour planifier les tournées optimales, r\'ealisant une r\'eduction de 20\% des co\^uts op\'erationnels et recon\'nue 1\textsuperscript{er} prix au CITS National (75 \'equipes). Cette exp\'erience couvre exactement les briques techniques que demande votre challenge : s\'eries temporelles, optimisation de routes, et scaling sur donn\'ees r\'eelles.

\vspace{0.8em}\noindent Au MILA (Montréal), j'ai conduit une recherche fondamentale rigoureuse sur la dynamique d'optimisation des r\'eseaux de neurones (phénomène grokking), combinant formulation math\'ematique, validation statistique et expériences `a grande \'echelle. Appliquée `a RTE, cette rigueur scientifique garantit que les mod\`eles de prédiction de charge ou planification de tournées r\'eseau seront auditable, conformes aux standards d'exploitation, et transférable d'un gestionnaire de r\'eseau `a l'autre.

\vspace{0.8em}\noindent Je maîtrise enfin une stack ML-ops moderne (Python, scikit-learn, PyTorch, Docker, CI/CD) et l'orchestration de syst\`emes `a base de LLM (Google ADK, LangChain) pour analyse documentaire (normes, donn\'ees historiques). Le syst\`eme Scoop repr\'esente un cas d'usage complexe : intégration de données r\'eseau temps r\'eel, optimisation multi-crit\`eres (t\'emps, \'energie, co\^uts), explainabilit\'e des décisions. RTE est pour moi un environnement id\'eal pour appliquer et renforcer ces comp\'etences.

\vspace{0.8em}\noindent Je suis disponible d\`es septembre 2026 pour une alternance de 12 mois. J'apporterai `a RTE une exp\'erience concrete en optimisation de tournées smart, fondamentaux scientifiques solides, et engagement `a d\'eveloper des syst\`emes data science qui soutiendront l'excellence op\'erationnelle du gestionnaire du r\'eseau \'electrique fran\c cais.

\vspace{1.5em}\noindent Veuillez agr\'eer, Madame, Monsieur, l'expression de mes salutations distingu\'ees.
\vspace{2em}\noindent \textbf{Lo\"ic Aron MBASSI EWOLO}
\end{document}
